\chapter{Gaia DR3 Wide Binaries: A Coefficient-Level Consistency Check}
\label{app:gaia}

This appendix records a comparison of Gaia DR3 wide-binary kinematics
with the external-field-corrected form of the coefficient-level
interpolator already examined on SPARC in Appendix~\ref{app:empirical}.
It is not a test of Theorems~\ref{thm:mT}--\ref{thm:Leff}, not a
derivation of \(\nu_{\sqrt{}}\), and not evidence that Problem~H is
solved.

\section{What is being compared}
\label{sec:gaia-what}

The relevant acceleration scale for a solar-neighbourhood binary is
not the internal Newtonian acceleration alone. The Galactic external
field at the Sun satisfies
\[
  g_{\mathrm{gal}}\approx(1.5\text{--}2.2)\times 10^{-10}\,\mathrm{m\,s^{-2}}
  >g_\dagger.
\]
Under the external-field saturation already noted for the
coefficient layer (Assumption~\ref{ass:efe}), the isolated
deep-regime boost is suppressed. The working prediction used here is
therefore the residual enhancement obtained by evaluating
\(\nu_{\sqrt{}}\) on the quadrature sum of the internal and external
fields (or the equivalent AQUAL/QUMOND external-field prescription).
Asymptotically this residual is of order \(20\%\), not the factor
\(\sim4\) of the isolated map.

The comparison is performed on the public El-Badry et al.\ Gaia DR3
wide-binary catalogue after the standard quality cuts
(RUWE\(<1.3\), fractional parallax error \(<0.06\), projected
separations \(2\)--\(30\,\mathrm{kAU}\)). Projection, unresolved
companions, and chance alignments are modelled by the same mixture
techniques employed in the controlled literature analyses; they are
not optional.

\section{What the data show}
\label{sec:gaia-show}

Controlled analyses that include the Galactic external field,
realistic triple contamination, and the quality cuts above find
relative velocities consistent with Newtonian gravity. The residual
MOND/RAR-like enhancement is excluded at high significance
(Banik et al.\ 2024; Pittordis, Sutherland \& Shepherd 2025;
quality-framework study 2026). Apparent boosts reported in less
filtered samples diminish once these systematics are treated.

Consequently the external-field-corrected \(\nu_{\sqrt{}}\) (the
coefficient-level map of Appendix~\ref{app:empirical}) is also
consistent with the data only to the extent that its residual
boost remains inside the observational uncertainties. It does not
outperform Newtonian gravity on this sample, nor does it produce a
statistically preferred rising velocity ratio of the size generated
by the isolated interpolator.

\section{What this is not}
\label{sec:gaia-not}

\begin{itemize}
  \item It is not a derivation of \(\nu_{\sqrt{}}\) from
        \(I(A:B)\) or from the modular spectrum. Problem~H remains
        open.
  \item It is not a constraint on \(m_T=0\), on the heat-kernel
        map for \(G_{\mathrm{eff}}\), or on the conservation
        identity. A null result under the external-field-corrected
        prediction is the outcome expected once
        Assumption~\ref{ass:efe} is respected.
  \item It is not evidence that an isolated reading of the RAR
        interpolator applies to solar-neighbourhood binaries.
        That reading is already ruled out by \(g_{\mathrm{gal}}>g_\dagger\).
  \item It is not a laboratory handle on the entanglement stress.
\end{itemize}

\section{Status}
\label{sec:gaia-status}

The Gaia DR3 wide-binary sample, analysed with the external field
included, is compatible with Newtonian gravity and with the
modest residual enhancement allowed by the coefficient-level
NAMCG/RAR map under external-field saturation. It supplies no
new support for the framework beyond the SPARC consistency check
already recorded, and it does not alter the open-problem ledger
of Appendix~\ref{app:status}.